% =============================================================================
% Spring force in the 11x11 lattice simulation — technical note
%
% Compile from repository root:
%   cd papers/technical && latexmk -pdf -output-directory=../build -aux-directory=../build spring_force_11x11.tex
% =============================================================================

\documentclass[11pt,a4paper]{article}
\usepackage[utf8]{inputenc}
\usepackage{amsmath,amssymb,amsfonts}
\usepackage{hyperref}
\usepackage[margin=1in]{geometry}
\usepackage{booktabs}
\usepackage{siunitx}
\usepackage{listings}
\usepackage{xcolor}
\usepackage{parskip}
\usepackage{microtype}
\usepackage{enumitem}
\usepackage{tcolorbox}
\tcbuselibrary{skins,breakable}

\setlist{itemsep=0.35em, topsep=0.35em}

% --- ``In code'' callout: theory block vs implementation ---
\newtcolorbox{codebox}{
  colback=gray!7,
  colframe=gray!55,
  boxrule=0.6pt,
  arc=2pt,
  left=8pt,
  right=8pt,
  top=6pt,
  bottom=6pt,
  breakable,
  title={\textbf{In the Julia code} \textnormal{(\texttt{src/lattice\_simulation\_11x11.jl})}},
  fonttitle=\small\bfseries,
  coltitle=black,
  attach boxed title to top left={yshift=-2mm, xshift=4mm},
  boxed title style={empty, frame hidden},
}

\lstset{
    basicstyle=\ttfamily\footnotesize,
    breaklines=true,
    frame=single,
    numbers=left,
    numberstyle=\tiny\color{gray},
    xleftmargin=1.5em,
    framexleftmargin=1em,
}

\title{\textbf{How the 11$\times$11 simulation computes spring forces}\\[0.35em]
\large Exponential springs, step by step (theory and code)}
\author{Daniel Miles}
\date{\today}

\begin{document}
\maketitle

\begin{abstract}
\noindent
This note explains the \emph{geometric} nonlinear spring model used in the blast-loaded sandwich-plate simulations, then shows \emph{where} each piece lives in the Julia program. The core is an exponential relation between bond stretch $s=|\mathbf{R}|-|\mathbf{R}_0|$ (current chord length minus reference length) and the axial force along the deformed chord. The lattice uses \textbf{nearest-neighbor bonds only} (horizontal and vertical); diagonal springs are not in the ODE or energy. Horizontal springs use global $(k,\alpha)$; vertical springs use column material $(k,\alpha,c)$. The potential energy routine matches the dynamics for the same bond routing.
\end{abstract}

% ---------------------------------------------------------------------------
\section{Why an exponential spring?}
\label{sec:why}

A linear spring (Hooke's law) says force is proportional to stretch. Real foams and many materials stiffen more strongly once deformations grow; an exponential law is one way to capture that in a mass--spring network without leaving the familiar picture of nodes connected by springs.

\medskip
\noindent\textbf{Plain-language picture.} Imagine two beads and a link between them. We pick a \emph{reference} layout (a regular grid). While the simulation runs, each bead's \textbf{displacement} is how far it has moved away from its reference spot. The program measures how much the two beads have shifted \textbf{relative to each other}; that relative shift sets the link's stretch proxy $d$. The push--pull along the link grows like $e^{\alpha d}-1$ (after a scale $k$), so small motions feel nearly linear, while larger ones ramp up faster.

\section{Reference lattice and what the solver actually stores}
\label{sec:kinematics}

\subsection{The undisturbed grid}
Think of an $11\times 11$ array of masses, rows $i$ and columns $j$, both running from $1$ to $N$ with $N=11$. In the reference (stress-free layout used only as a bookkeeping frame), the position of mass $(i,j)$ is
\begin{equation}
\mathbf{X}_{ij}^{\mathrm{eq}} =
\begin{bmatrix} (j-1)\,\Delta \\ (i-1)\,\Delta \end{bmatrix},
\qquad \Delta = 1.0\ \mathrm{m}.
\end{equation}
So neighbors along a row sit $\Delta$ apart horizontally; neighbors along a column sit $\Delta$ apart vertically.

\subsection{Displacements versus ``true'' positions}
The differential equation solver does \emph{not} store full positions $\mathbf{X}_{ij}(t)$ in the half of the state vector devoted to ``positions.'' It stores \textbf{displacements} $\mathbf{u}_{ij}(t)$ from that reference:
\begin{equation}
\mathbf{X}_{ij}(t) = \mathbf{X}_{ij}^{\mathrm{eq}} + \mathbf{u}_{ij}(t).
\end{equation}
At the start of a run, every mass is at rest at its reference node: $\mathbf{u}_{ij}(0)=\mathbf{0}$. When you see animations, the plotted positions add the reference grid back in.

\begin{codebox}
The reference grid is \texttt{create\_equilibrium\_grid()} using \texttt{GRID\_SPACING = 1.0}. The state vector \texttt{u} begins as \texttt{u0 = zeros(...)}; its position block is reshaped into \texttt{pos}, with \texttt{pos[:,k] = u}$_{ij}$ for index \texttt{k = lattice\_idx(i,j)}. Plots use \texttt{current = displacement + equilibrium\_grid}.
\end{codebox}

\subsection{Reference chord and deformed chord}
For a bond from mass $a$ to neighbor $b$, let $\mathbf{R}_0=\mathbf{X}_b^{\mathrm{eq}}-\mathbf{X}_a^{\mathrm{eq}}$ (reference chord from $a$ toward $b$). With stored displacements,
\begin{equation}
\mathbf{R}=\mathbf{R}_0+(\mathbf{u}_b-\mathbf{u}_a),\qquad
L_0=\|\mathbf{R}_0\|,\qquad L=\|\mathbf{R}\|,\qquad s=L-L_0.
\label{eq:rdef}
\end{equation}
At $\mathbf{u}=0$, $\mathbf{R}=\mathbf{R}_0$ and $s=0$: \textbf{no spring force}. The stretch $s$ is the true elongation of the bond in the lab frame (geometric spring), and $\hat{\mathbf{R}}=\mathbf{R}/L$ is the axial direction under shear or combined motion.

% ---------------------------------------------------------------------------
\section{Force on one spring: magnitude and direction}
\label{sec:force-law}

With $\mathbf{R}$, $L$, and $s$ as in \eqref{eq:rdef}, if $L$ is above a tiny floor ($10^{-12}$ in code),
\begin{align}
\hat{\mathbf{R}} &= \frac{\mathbf{R}}{L}, \\
F_{\mathrm{mag}} &= k\left(e^{\alpha s}-1\right), \\
\mathbf{F}_{\text{on }a\,\text{from }b} &= F_{\mathrm{mag}}\,\hat{\mathbf{R}}.
\label{eq:force}
\end{align}
If $L$ is below the floor, the force is set to zero. The direction follows the \emph{deformed} chord $\mathbf{R}$, so shear rotates the axial force correctly.

\medskip
\noindent\textbf{Units.} $k$ is in N, $\alpha$ in m$^{-1}$, and $s$ in m, so $F_{\mathrm{mag}}$ is a force. For small $|s|$, $F_{\mathrm{mag}}\approx k\alpha s$.

\begin{codebox}
One bond is evaluated in \texttt{spring\_force\_2d\_geometric(uA, uB, R0, k, alpha)}. The RHS \texttt{lattice\_2d\_rhs\_nn\_backplate!} passes the correct $\mathbf{R}_0$ for left/right/up/down neighbors (scaled by \texttt{GRID\_SPACING}).
\end{codebox}

% ---------------------------------------------------------------------------
\section{Stored energy and the nearly-linear regime}
\label{sec:energy}

The potential energy per bond consistent with \eqref{eq:force} along the bond is
\begin{equation}
U(s) = \frac{k}{\alpha}\left(e^{\alpha s} - \alpha s - 1\right),
\qquad s=L-L_0.
\end{equation}
For small $|s|$, $U\approx \tfrac{1}{2}k\alpha\,s^2$.

\begin{codebox}
\texttt{potential\_energy\_2d\_nn\_backplate} sums this $U(s)$ over all NN bonds with the same $(k,\alpha)$ routing as \texttt{lattice\_2d\_rhs\_nn\_backplate!}, plus the wall penalty.
\end{codebox}

% ---------------------------------------------------------------------------
\section{Wiring the whole lattice: neighbors, layers, and dashpots}

\subsection{Who is connected to whom?}
Each interior mass has up to four NN neighbors (left, right, up, down). Edge and corner nodes have fewer. There are no diagonal springs in the 11×11 simulation.

\subsection{Which springs are ``softer'' or ``stiffer'' by design?}
\begin{itemize}
  \item \textbf{Horizontal links.}\enspace Global $(K_{\text{COUPLING}},\alpha_{\text{COUPLING}})$ and base $C_{\text{DAMPING}}$ (not column material).
  \item \textbf{Vertical links} (same column, adjacent rows) use column $j$ material: $k$, $\alpha$, and $c$ from \texttt{get\_column\_*}.
  \item \textbf{Dashpots} act on all four NN directions, with column-based $c$ on vertical bonds.
  \item Tuple fields \texttt{k\_diagonal}, \texttt{alpha\_diagonal} exist for optimizer/API compatibility but are \textbf{unused} in the dynamics.
\end{itemize}

\begin{table}[htbp]
\centering
\small
\begin{tabular}{@{}p{0.22\textwidth}p{0.43\textwidth}p{0.27\textwidth}@{}}
\toprule
\textbf{Bond} & \textbf{Stiffness $(k,\alpha)$} & \textbf{Damping} \\
\midrule
Horizontal & $K_{\text{COUPLING}}$, $\alpha_{\text{COUPLING}}$ everywhere & $C_{\text{DAMPING}}$ \\
Vertical & Column $j$: \texttt{get\_column\_k\_coupling}, \texttt{get\_column\_alpha\_coupling} & \texttt{get\_column\_c\_damping} \\
\bottomrule
\end{tabular}
\caption{NN bond parameters in \texttt{lattice\_2d\_rhs\_nn\_backplate!}.}
\label{tab:routing}
\end{table}

\subsection{Default material scaling}
The program builds $N=11$ material records. Material $1$ uses the base numbers; for index $m>1$, stiffnesses and damping scale by $\sigma^{m-1}$ with $\sigma=\texttt{MATERIAL\_MULTIPLIER}$ (default $2/3$), while $\alpha$ is usually kept at the base value. A \texttt{material\_order} permutation places those materials in columns for optimization studies.

\begin{codebox}
Material tuples are built in \texttt{create\_default\_material\_properties}. During integration, parameters come from \texttt{p.material\_order} and \texttt{p.materials} when present; otherwise default column scaling uses \texttt{DEFAULT\_MATERIALS} with identity column order. The ODE right-hand side is \texttt{lattice\_2d\_rhs\_nn\_backplate!}.
\end{codebox}

\subsection{Viscous damping (nearest neighbors only)}
Between nearest neighbors, an extra piece proportional to relative velocity is added on mass $a$:
\begin{equation}
\mathbf{F}_{\mathrm{damp}} = -c\left(\mathbf{v}_a-\mathbf{v}_b\right),
\end{equation}
with $c$ taken from the same routing rules as the vertical spring for column neighbors (or globals for horizontal neighbors).

% ---------------------------------------------------------------------------
\section{Forces that are not springs}
These are separate from the exponential links but appear in the same acceleration routine:
\begin{itemize}
  \item \textbf{Backplate.} Rightmost column: if absolute $x$ passes the plate, a stiff spring-like penalty (and damping on $v_x$ pushing into the wall) pushes back.
  \item \textbf{Blast-type load.} For a short active window, trapezoidally distributed forces drive the left edge (column~1).
\end{itemize}

\begin{codebox}
Wall logic uses absolute $x = \texttt{(N-1)*GRID\_SPACING + pos\_k[1]}$. Loads use \texttt{F\_ACTIVE\_TIME}, \texttt{F\_MAG}, and row weights \texttt{get\_distributed\_force\_magnitude}.
\end{codebox}

% ---------------------------------------------------------------------------
\section{Default numerical values in the simulation}
\label{sec:constants}

Table~\ref{tab:constants} lists the global constants as in the current source file. Vertical NN springs use column materials; diagonal constants in the table are retained in tuples only and are not applied in the ODE.

\begin{table}[htbp]
\centering
\small
\begin{tabular}{@{}p{0.36\textwidth}p{0.18\textwidth}p{0.38\textwidth}@{}}
\toprule
\textbf{Name / symbol} & \textbf{Value} & \textbf{Role} \\
\midrule
$N$ & $11$ & Grid side length \\
$\Delta$ (\texttt{GRID\_SPACING}) & \SI{1.0}{m} & Reference spacing \\
$m$ (\texttt{MASS}) & \SI{1.0}{kg} & Mass per node \\
$K_{\text{COUPLING}}$ & \SI{100}{N} & Horizontal NN springs; vertical scale reference \\
$\alpha_{\text{COUPLING}}$ & \SI{10}{m^{-1}} & NN exponential rate (code comments say ``decay'' but here $\alpha>0$ \emph{stiffens} with stretch) \\
$K_{\text{DIAGONAL}}$ & \SI{50}{N} & Unused in ODE (kept in material tuples) \\
$\alpha_{\text{DIAGONAL}}$ & \SI{10}{m^{-1}} & Unused in ODE (kept in material tuples) \\
$C_{\text{DAMPING}}$ & \SI{5}{N\,s\,m^{-1}} & Base NN dashpot \\
$\sigma$ (\texttt{MATERIAL\_MULTIPLIER}) & $2/3$ & Softens $k,c$ for higher material indices \\
$k_{\mathrm{wall}}$, $c_{\mathrm{wall}}$ & \SI{10000}{N\,m^{-1}}, \SI{10}{N\,s\,m^{-1}} & Backplate \\
$d_{\text{back}}$ & \SI{0.001}{m} & Offset to plate \\
$F_{\mathrm{tot}}$, $t_{\mathrm{act}}$ & \SI{1200}{N}, \SI{0.15}{s} & Edge load budget and duration \\
\bottomrule
\end{tabular}
\caption{Global parameters (\texttt{const} values in \texttt{src/lattice\_simulation\_11x11.jl}).}
\label{tab:constants}
\end{table}

% ---------------------------------------------------------------------------
\section{Matching symbols in the written paper to this program}
\label{sec:notation}

The manuscript writes spring forces using something like $\mathbf{r} = \mathbf{x}_{i',j'}-\mathbf{x}_{i,j}$. In the code, $\mathbf{x}$ must mean the \textbf{same displacements} stored in \texttt{pos}, not absolute lab coordinates $\mathbf{X}$. Starting from $\mathbf{u}=\mathbf{0}$ is the same as starting from $\mathbf{X}=\mathbf{X}^{\mathrm{eq}}$ in the paper. If you accidentally plugged absolute positions into the same formula, every bond would appear stretched at $t=0$ and the model would predict huge initial forces, which is not what the program does.

% ---------------------------------------------------------------------------
\section*{Appendix: core routine (Julia)}
\begin{lstlisting}
function spring_force_2d_geometric(uA, uB, R0, k, alpha)
    R = R0 .+ (uB .- uA)
    L = norm(R)
    if L < 1e-12
        return zeros(2)
    end
    L0 = norm(R0)
    s = L - L0
    Fmag = k * (exp(alpha * s) - 1.0)
    return Fmag * (R / L)
end
\end{lstlisting}

\end{document}
